% Generated draft: requires mathematical and source review.
\documentclass[11pt]{article}
\usepackage[margin=25mm]{geometry}
\usepackage{fontspec}
\setmainfont{Latin Modern Roman}
\usepackage{amsmath,amssymb,mathtools,unicode-math}
\setmathfont{Latin Modern Math}
\usepackage{enumitem,needspace,xurl,hyperref}
\setlength{\parindent}{0pt}
\setlength{\parskip}{5pt}
\setlength{\emergencystretch}{3em}
\newcommand{\N}{\mathbb{N}}
\newcommand{\Z}{\mathbb{Z}}
\newcommand{\Q}{\mathbb{Q}}
\newcommand{\R}{\mathbb{R}}
\newcommand{\C}{\mathbb{C}}
\newcommand{\F}{\mathbb{F}}
\newcommand{\A}{\mathbb{A}}
\newcommand{\PP}{\mathbb{P}}
\newcommand{\E}{\mathbb{E}}
\newcommand{\eps}{\varepsilon}
\newcommand{\dd}{\,\mathrm{d}}
\begin{document}
\section*{UnsolvedMath ID: 6200083}
\textbf{Source problem code:} catalogue entry
\subsection{Definitions and mathematical statement}
Let \(X\) be a Gromov-hyperbolic space with its intended Gromov boundary \(\partial X\), and let \(G\le\operatorname{Isom}(X)\) be finitely generated and discrete. For a base point \(x_0\in X\), the limit set is the boundary accumulation set
\[
\Lambda_X(G)=\{\xi\in\partial X:\exists(g_j)\subseteq G,
\ g_jx_0\longrightarrow\xi\text{ in }X\cup\partial X\}.
\]
Use the subspace topology from the Gromov boundary. Local connectedness means that for every \(\xi\in\Lambda_X(G)\) and every relative open neighborhood \(U\) of \(\xi\), there is a connected relative open set \(V\) with \(\xi\in V\subseteq U\).

Kapovich's question is the implication
\[
\forall(X,G)\text{ as above},\qquad
\Lambda_X(G)\text{ connected}\quad\Longrightarrow\quad
\Lambda_X(G)\text{ locally connected}.
\]
No cocompactness, quasiconvexity, or Kleinian/arithmetic hypothesis is stated. In particular, the limit set is not silently identified with the abstract boundary of \(G\), since the hypotheses do not say that \(G\) is itself hyperbolic or that an orbit map is a quasi-isometric embedding.

The extraction does not specify properness or geodesicity of \(X\), a compactification convention in the nonproper case, or whether discrete means a discrete subgroup in the isometry-group topology or properly discontinuous action. In the standard proper geodesic setting the boundary and orbit-limit definitions are unambiguous; restricting to that setting would be an additional hypothesis requiring confirmation. The supplied research summary reports no general proof and mentions Cannon--Thurston/local-connectivity work, but supplies no precise theorem for this full class. Its vague non-arithmetic/standard-hypothesis remark is not made an assumption.
\subsection{Short English statement}
Must a connected limit set of a finitely generated discrete isometry group of a Gromov-hyperbolic space be locally connected? Preserve the broad space class unless properness or other missing action conventions are confirmed.
\subsection{Sources}
\begin{itemize}
\item Local UnsolvedMath registry, permanent ID \(6200083\), exact source code \texttt{catalogue entry}, ``Boundaries of Groups and Kleinian Groups --- Problem 83.'' Proposer Misha Kapovich, 2005.
\item Kapovich, \emph{Problems on Boundaries of Groups and Kleinian Groups} (2005), Problem 83, recorded PDF page 21; \url{https://www.math.ucdavis.edu/~kapovich/EPR/problems.pdf}. Recorded access 2026-07-29; pdf-text-normalized extraction.
\item Supplied OPEN-TRIAGE research summary invokes a general Cannon--Thurston/local-connectivity programme without complete citations or exact hypotheses for any comparison theorem.
\end{itemize}
\end{document}
